%% Formule de l'operateur : UV2ddff
%% Dans ce fichier, les �ventuels accents sont cod�s
%% en ISO-8859-15 (aka latin9)
%% La ligne suivante donne trois exemples 
%% e accent aigu : �, e accent grave : �, e-dans-l'o : �

$FF = \sqrt{ U^2 + V^2 } $ \newline

Si $V=0$ alors  \newline
\begin{tabular}{p{3cm}|l}
 & si $U=0$ alors $DD=0$ \\
 & si $U > 0$ alors $DD=-\pi /2$ \\
 & si $U < 0$ alors $DD=\pi /2$ \\
\end{tabular}

Si $V > 0$ alors \newline
\begin{tabular}{p{3cm}|l}
 & $DD = \arctan ( U / V)  + \pi$ \\
\end{tabular}

Si $V < 0$ alors \newline
\begin{tabular}{p{3cm}|l}
 & $DD = \arctan ( U / V)$  \\
\end{tabular}
\newline
Les �ventuelles directions n�gatives sont rendues positives : si $DD < 0$ alors
$ DD = DD + 2*\pi $

Puis conversion en degr�s :
$ DD = DD * \frac{180}{\pi}$

